\chapter{Conclusion}

\section{Main Findings}

The results are interpreted as paired comparisons over the same elapsed
time rather than as rankings based on unmatched terminal values. Formal
simulations use a pre-specified 1000-business-day horizon and stop early
only when one or fewer non-central banks remain active. Runs that reach
day 1000 without collapse are right-censored for collapse-time analysis.

Under the baseline configuration, in which rollover and policy support
are both enabled, average systemic risk remains low under both matching
mechanisms through day 1000. The paired full-horizon mean is 0.0396 for
decentralized RFQ and 0.0447 for centralized matching. The mean paired
RFQ-minus-Centralized difference is \(-0.00513\), with a 95\% interval
of \([-0.01276,0.00249]\). The point estimate therefore favors RFQ, but
the interval narrowly includes zero and does not establish separation at
the conventional 5\% level.

The endpoint comparison points in the same direction: the paired
day-1000 difference is \(-0.01727\), with interval
\([-0.03233,-0.00222]\). The day-1000 CGR difference is
\(-0.07555\), whose interval \([-0.12626,-0.02484]\) excludes zero.
Thus the strongest baseline evidence concerns lower accumulated capital
shortfall under RFQ, while the composite-risk result is directionally
favorable but statistically less decisive.

Policy conclusions are based on risk paths and bank-days rather than
collapse dates alone because the baseline ON/ON sample is fully
right-censored at day 1000. Relative to OFF/OFF, ON/ON increases average
non-central-bank days alive by 13,353.9 in RFQ and 16,004.1 in centralized
matching on a fixed 1000-day horizon; both paired intervals are positive. This supports a
stabilizing interpretation of the combined intervention, while the
complete censoring of collapse time prevents claims about mean time to
systemic collapse beyond the observed horizon.

Rollover should be interpreted as a contractual maturity treatment. With
rollover ON, new loans are originated as 5--20-period installment
contracts and unpaid installments become arrears; with rollover OFF, new
loans mature as next-day single payments and any residual after clearing
causes hard default. The four-scenario comparison therefore identifies
the interaction between maturity design and policy support, not a
mechanical ex-post extension of the same bullet contract.

The matching rule has a stronger and more precisely estimated effect on
network structure than on the composite-risk level. RFQ creates more
active links (58.95 versus 37.77) and a larger connected component (24.03
versus 17.08), but lower exposure HHI (0.050 versus 0.094) and a much
smaller largest bilateral exposure (116 versus 1,119). RFQ also has lower
funding satisfaction (0.530 versus 0.597), while cumulative volume is not
statistically separated. The economic interpretation is therefore not
that RFQ supplies more funds; it distributes activity across a broader
and less concentrated set of bilateral claims.

Both matching mechanisms operate every business day. Centralized
matching observes the full eligible market and commits allocations only
when a borrower's complete demand can be covered. RFQ searches a bounded
local information pool, applies GNN-assisted screening, and retains
acceptable partial funding. The estimated difference therefore reflects
information scope, screening, allocation, pricing, and partial-funding
rules, but is not driven by unequal trading frequency.

The weight sweep preserves the direction of the average RFQ advantage
over the tested \(w_1\) grid, but the CAR-threshold sweep is not uniformly
one-sided. RFQ has lower measured SR at the baseline threshold 0.08 and
at 0.09, while the ranking reverses from 0.10 onward in
the current sample. This is evidence that the substantive conclusion is
sensitive to how broadly capital stress is classified and reinforces the
need to report \(\mathrm{FR}_t\), \(\mathrm{CBS}_t\), and
\(\mathrm{CGR}_t\) separately from the composite measure.

\section{Contributions}

This thesis makes four main contributions.

First, it develops a dynamic interbank network framework in which
bilateral lending relationships are generated endogenously from changing
bank states and trading decisions rather than imposed as a fixed network
before contagion is evaluated.

Second, it implements and compares two daily transaction designs within
the same balance-sheet, regulatory, contract-accounting, and clearing
framework: a market-wide centralized mechanism with an all-or-nothing
borrower rule and a decentralized RFQ mechanism with bounded local lender
discovery, bilateral quotation, GNN-assisted screening, and retained
partial funding. This makes market design itself part of the endogenous
network-formation process.

Third, it introduces a dated contract-book structure that separates the
stock of outstanding exposure from cash flows currently due. The
rollover switch selects either a 5--20-period amortizing contract or a
next-day single-payment contract at origination. Only scheduled cash flows enter the gross Eisenberg--Noe
liability matrix, allowing repayment timing and creditor losses to feed
back into later bank states without prematurely clearing future
obligations.

Fourth, it combines multidimensional risk measurement with a
right-censoring-aware comparison over a pre-specified horizon. The
framework records failures, capital stress, capital gaps, network
structure, funding outcomes, project losses, and policy intervention.
Paired state paths, bank-days, and collapse indicators are kept distinct,
which is important when an intervention reduces vulnerability without
producing an observed systemic collapse within the study horizon.

Overall, the contribution of the thesis lies in connecting the trading
process that creates interbank exposures with the contractual timing and
clearing process through which those exposures later generate funding
pressure and systemic loss.

\section{Limitations and Future Research}

The numerical results should be interpreted within the scope of the
stylized simulation specified in Chapter~6. The system contains 30 institutions whose balance
sheets, funding structures, behavioral parameters, and market conditions
are generated from model-based rules rather than calibrated to
institution-level empirical data. Monetary values are simulation units,
and the reported systemic-risk levels are comparative model outcomes
rather than forecasts for a particular banking system.

The matching mechanisms simplify actual interbank trading. Both operate
daily, but the centralized benchmark combines market-wide information
with an all-or-nothing funding rule, while RFQ combines local discovery,
learned screening, and partial execution. The experiment therefore
identifies the difference between these complete transaction designs,
not the isolated causal effect of any single component. Future work
could separately randomize information scope, screening, pricing, and
partial-funding permission.

The rollover mechanism is also stylized. It assigns an amortizing
5--20-period schedule at origination without modeling a full refinancing
market, collateral negotiation, endogenous lender consent, or market
pricing of maturity extension. Policy support uses
rule-based liquidity and solvency/capital facilities with finite budgets
rather than a complete institutional representation of central-bank
operations.

The 1000-day horizon is long relative to the daily contract cycle. The
baseline ON/ON sample is fully right-censored at that horizon, so
collapse time cannot rank mechanisms there. Intermediate first-passage
thresholds for \(\mathrm{SR}_t\) are reported in Chapter~7; high
failure-rate thresholds are not reached in ON/ON. OFF/OFF centralized
matching does record early stops, which shows that collapse is identified
only in some contractual environments. Longer simulations, larger Monte
Carlo samples, and event-history methods for additional intermediate
thresholds would provide more information about tail stability.

The contagion analysis focuses on contractual interbank payments and
balance-sheet feedback. Other channels, including common asset holdings,
fire-sale price effects, deposit runs, margin calls, collateral spirals,
and interactions with non-bank sectors, remain outside the current
framework. The capital and liquidity indicators are model-based proxies
and do not constitute a complete implementation of all Basel III
regulatory calculations.

Future research can calibrate balance sheets, maturities, bilateral
limits, quotation behavior, policy parameters, and shock processes with
empirical data. A larger banking population and richer network structure
would permit stronger tests of heterogeneity. Structural robustness
could also be expanded beyond the bilateral loan cap to include network
size, bank-type composition, RFQ search depth, maximum RFQ rounds,
minimum trade size, centralized auction frequency, maturity structure,
liquidity buffers, policy budgets, and shock intensity.
